考虑三能级 \(\Lambda\) 型体系的正交基 \(|1\rangle,|2\rangle,|3\rangle\)，在旋转波近似与双光子共振条件下，其旋转参考系哈密顿量为
\[
H(t)=\frac{\hbar}{2}
\begin{pmatrix}0&\Omega_p(t)&0\\
\Omega_p(t)&2\Delta&\Omega_s(t)\\
0&\Omega_s(t)&0\end{pmatrix},
\]
其中 \(\Delta\) 为固定实单光子失谐，\(\Omega_p,\Omega_s\) 为可控、实且非负的 Rabi 频率；脉冲交叠时 \(\Omega=\sqrt{\Omega_p^2+\Omega_s^2}>0\)。
(1) 严格对角化瞬时哈密顿量，求零能暗态、两个亮态本征能量以及暗态中间能级的占据概率。
(2) 定义混合角 \(\tan\theta=\Omega_p/\Omega_s\)，求暗态随时间变化产生的非绝热耦合，并推导相对最近亮态的最小能隙；给出明确的绝热充分量级条件。
(3) 对先开 Stokes 脉冲、后开 pump 脉冲的反直觉序列，证明理想绝热跟随可把 \(|1\rangle\) 转移到 \(-|3\rangle\)，且不占据易衰变的 \(|2\rangle\)。讨论实耦合时暗态 Berry 联络。
(4) 加入小双光子失谐扰动 \(V=\hbar\delta|3\rangle\langle3|\)。计算暗态的一阶能移及其与明态基矢的耦合，说明失谐为何破坏严格暗态；再求 \(|\Delta|\gg\Omega\) 时最近亮态间隙的主导渐近，并讨论对绝热制备的影响。
